% GENERATED FILE. Edit canonical lessons, then run npm run generate:books.
% canonical-source-hash: fnv1a64:873c9f7f2e61c470
% canonical-lessons: TE-C211-timmiri, TE-C211-ekkillu, TE-C211-guraka, TE-C211-vadadebba, TE-C211-aviri

\chapter{Numbness, Hiccups, Snoring, Sunstroke, Steam}
\label{ch:te-numbness-hiccups-snoring-sunstroke-steam}
\hlchaptermodality{211}

\begin{chapteropening}
\textbf{By the end of this chapter:} \emph{I can say numbness, hiccups, snoring, sunstroke and steam.}

Complete the last lesson of chapter 211: I can say numbness, hiccups, snoring, sunstroke and steam.
\end{chapteropening}

\section[\texorpdfstring{\te{తిమ్మిరి} (timmiri)}{తిమ్మిరి (timmiri)}]{\te{తిమ్మిరి} (timmiri) — numbness, pins and needles}

\label{lesson:TE-C211-timmiri}



\begin{quote}
Before the new one: say the Telugu for boredom, then the Telugu for pining.
\end{quote}

\subsection*{You'll want to know: \te{తిమ్మిరి}}
\textbf{\te{తిమ్మిరి}} — \emph{timmiri} — \textquotedblleft{}numbness, pins and needles\textquotedblright{}.

\begin{cousinweb}
\textbf{\te{తిమ్మిరి}} (\emph{timmiri}) is native Dravidian. DEDR 3235 lists \emph{timmiri}, \textquotedblleft{}numbness\textquotedblright{}, with Tamil \emph{timir} and Malayalam \emph{timir}, \textquotedblleft{}numbness\textquotedblright{}.
\end{cousinweb}

\begin{grammarlens}[title={What you've built}]
One more word for this chapter.
\end{grammarlens}

\subsection*{Your turn}
\begin{itemize}
\raggedright
  \item \emph{Say it:} \emph{timmiri}
  \item \emph{Say it:} \emph{timmiri}, once more
  \item \emph{Say it:} say \emph{timmiri}
  \item \emph{Recall:} say the Telugu for boredom, then the Telugu for pining, then say \emph{timmiri} again
\end{itemize}

\begin{tcolorbox}[breakable,colback=teal!4,colframe=teal!35!black,title={Before you move on}]
What does \te{తిమ్మిరి} mean? (\textquotedblleft{}Numbness, pins and needles\textquotedblright{}.) Say it once more.
\end{tcolorbox}
\section[\texorpdfstring{\te{ఎక్కిళ్ళు} (ekkiḷḷu)}{ఎక్కిళ్ళు (ekkiḷḷu)}]{\te{ఎక్కిళ్ళు} (ekkiḷḷu) — hiccups}

\label{lesson:TE-C211-ekkillu}



\begin{quote}
Before the new one: say the Telugu for pining, then the Telugu for numbness.
\end{quote}

\subsection*{You'll want to know: \te{ఎక్కిళ్ళు}}
\textbf{\te{ఎక్కిళ్ళు}} — \emph{ekkiḷḷu} — \textquotedblleft{}hiccups\textquotedblright{}.

\begin{cousinweb}
\textbf{\te{ఎక్కిళ్ళు}} (\emph{ekkiḷḷu}) is native Dravidian. DEDR 772 lists \emph{ekkili}, \textquotedblleft{}hiccup\textquotedblright{}, with Malayalam \emph{ekkiḷ}, \textquotedblleft{}hiccough\textquotedblright{}; \textbf{\te{ఎక్కిళ్ళు}} (\emph{ekkiḷḷu}) is its plural.
\end{cousinweb}

\begin{grammarlens}[title={What you've built}]
One more word for this chapter.
\end{grammarlens}

\subsection*{Your turn}
\begin{itemize}
\raggedright
  \item \emph{Say it:} \emph{ekkiḷḷu}
  \item \emph{Say it:} \emph{ekkiḷḷu}, once more
  \item \emph{Say it:} say \emph{ekkiḷḷu}
  \item \emph{Recall:} say the Telugu for pining, then the Telugu for numbness, then say \emph{ekkiḷḷu} again
\end{itemize}

\begin{tcolorbox}[breakable,colback=teal!4,colframe=teal!35!black,title={Before you move on}]
What does \te{ఎక్కిళ్ళు} mean? (\textquotedblleft{}Hiccups\textquotedblright{}.) Say it once more.
\end{tcolorbox}
\section[\texorpdfstring{\te{గురక} (guraka)}{గురక (guraka)}]{\te{గురక} (guraka) — snoring}

\label{lesson:TE-C211-guraka}



\begin{quote}
Before the new one: say the Telugu for numbness, then the Telugu for hiccups.
\end{quote}

\subsection*{You'll want to know: \te{గురక}}
\textbf{\te{గురక}} — \emph{guraka} — \textquotedblleft{}snoring\textquotedblright{}.

\begin{cousinweb}
\textbf{\te{గురక}} (\emph{guraka}) is native Dravidian. DEDR 1852 lists \emph{guṟaka}, \textquotedblleft{}snoring\textquotedblright{}, with Kannada \emph{guṟake} and Tamil \emph{kuṟaṭṭai}, \textquotedblleft{}snoring\textquotedblright{}, and points to related sound-words across the region (CDIAL 4207).
\end{cousinweb}

\begin{grammarlens}[title={What you've built}]
One more word for this chapter.
\end{grammarlens}

\subsection*{Your turn}
\begin{itemize}
\raggedright
  \item \emph{Say it:} \emph{guraka}
  \item \emph{Say it:} \emph{guraka}, once more
  \item \emph{Say it:} say \emph{guraka}
  \item \emph{Recall:} say the Telugu for numbness, then the Telugu for hiccups, then say \emph{guraka} again
\end{itemize}

\begin{tcolorbox}[breakable,colback=teal!4,colframe=teal!35!black,title={Before you move on}]
What does \te{గురక} mean? (\textquotedblleft{}Snoring\textquotedblright{}.) Say it once more.
\end{tcolorbox}
\section[\texorpdfstring{\te{వడదెబ్బ} (vaḍadebba)}{వడదెబ్బ (vaḍadebba)}]{\te{వడదెబ్బ} (vaḍadebba) — sunstroke}

\label{lesson:TE-C211-vadadebba}



\begin{quote}
Before the new one: say the Telugu for hiccups, then the Telugu for snoring.
\end{quote}

\subsection*{You'll want to know: \te{వడదెబ్బ}}
\textbf{\te{వడదెబ్బ}} — \emph{vaḍadebba} — \textquotedblleft{}sunstroke\textquotedblright{}.

\begin{cousinweb}
\textbf{\te{వడదెబ్బ}} (\emph{vaḍadebba}) is native Dravidian: \emph{vaḍa}, \textquotedblleft{}hot wind, heat\textquotedblright{}, one word with Kannada \emph{vaḍi}, \textquotedblleft{}heat\textquotedblright{} (DEDR 5225), and \emph{debba}, \textquotedblleft{}blow\textquotedblright{}, one word with Kannada \emph{debbe} (DEDR 3075).
\end{cousinweb}

\begin{grammarlens}[title={What you've built}]
One more word for this chapter.
\end{grammarlens}

\subsection*{Your turn}
\begin{itemize}
\raggedright
  \item \emph{Say it:} \emph{vaḍadebba}
  \item \emph{Say it:} \emph{vaḍadebba}, once more
  \item \emph{Say it:} say \emph{vaḍadebba}
  \item \emph{Recall:} say the Telugu for hiccups, then the Telugu for snoring, then say \emph{vaḍadebba} again
\end{itemize}

\begin{tcolorbox}[breakable,colback=teal!4,colframe=teal!35!black,title={Before you move on}]
What does \te{వడదెబ్బ} mean? (\textquotedblleft{}Sunstroke\textquotedblright{}.) Say it once more.
\end{tcolorbox}
\section[\texorpdfstring{\te{ఆవిరి} (āviri)}{ఆవిరి (āviri)}]{\te{ఆవిరి} (āviri) — steam, vapour}

\label{lesson:TE-C211-aviri}



\begin{quote}
Before the new one: say the Telugu for snoring, then the Telugu for sunstroke.
\end{quote}

\subsection*{You'll want to know: \te{ఆవిరి}}
\textbf{\te{ఆవిరి}} — \emph{āviri} — \textquotedblleft{}steam, vapour\textquotedblright{}.

\begin{cousinweb}
\textbf{\te{ఆవిరి}} (\emph{āviri}) is native Dravidian. DEDR 393 lists \emph{āviri}, \textquotedblleft{}vapour\textquotedblright{}, with Tamil \emph{āvi}, \textquotedblleft{}steam, vapour\textquotedblright{}, and Kannada \emph{āviri}, \textquotedblleft{}steam\textquotedblright{}.
\end{cousinweb}

\begin{grammarlens}[title={What you've built}]
One more word for this chapter.
\end{grammarlens}

\subsection*{Your turn}
\begin{itemize}
\raggedright
  \item \emph{Say it:} \emph{āviri}
  \item \emph{Say it:} \emph{āviri}, once more
  \item \emph{Say it:} say \emph{āviri}
  \item \emph{Recall:} say the Telugu for snoring, then the Telugu for sunstroke, then say \emph{āviri} again
\end{itemize}

\begin{tcolorbox}[breakable,colback=teal!4,colframe=teal!35!black,title={Before you move on}]
What does \te{ఆవిరి} mean? (\textquotedblleft{}Steam, vapour\textquotedblright{}.) Say it once more.
\end{tcolorbox}
